\documentclass{article}

\usepackage{amsmath,amssymb}
\usepackage{braket}
\usepackage{xurl}
\usepackage[colorlinks=true,linkcolor=blue,citecolor=blue,urlcolor=blue]{hyperref}
\setlength{\emergencystretch}{2em}

\input{command}

\newcommand{\ABAB}{\textit{ABAB}}
\newcommand{\ABA}{\textit{ABA}}
\newcommand{\BAB}{\textit{BAB}}
\newcommand{\closenessOfIP}{\textsf{closenessOfIP}}

\title{The evaluated-slice commutativity estimate for the
processed-$G$ measurements}
\author{}
\date{\today}

\begin{document}
\maketitle

\section{The assertion}

We examine the evaluated-slice commutativity lemma for the
processed $G$ measurements\footnote{The lemma carries the label
\leanid{lem:comm-data-processed-g} in
\path{references/ldt-paper/commutativity-G.tex}, lines~16--47, and is
tracked in \ghissue{760}.}
in \cite{Ji2020LowIndividualDegree}. Under the three hypotheses of
consistency with $A$, strong self-consistency, and boundedness, the
$G$ measurements approximately commute after evaluation for distinct
point pairs. The conclusion is
\begin{equation}\label{eq:processed-conclusion}
  \E_{\bu,\bv,\bx,\by} \sum_{a,b}
    \bigl\|\bigl(G^{\bx}_{[g(\bu)=a]} G^{\by}_{[h(\bv)=b]}
      - G^{\by}_{[h(\bv)=b]} G^{\bx}_{[g(\bu)=a]}\bigr)
    \ot I \ket{\psi}\bigr\|^{2}
  \leq 
  \nu, \qquad
  \nu = 48m\,(\sqrt{\gamma} + \sqrt{\zeta}).
\end{equation}

\section{The chain of approximations}

The paper's proof expands the commutator square and isolates the
$\ABAB$ term. The path from $\E[\sum \ABAB]$ to $\E[\sum \ABA]$
passes through ten approximation steps:
\begin{enumerate}
  \item Apply $\closenessOfIP$ with $A = G^{\bu,\bx}_a \ot I$ and
    $B = I \ot A^{\bu,\bx}_a$, at cost $2\sqrt{\zeta}$.
  \item Swap $A^{\bv,\by}_b$ onto the right register, at cost
    $\sqrt{\zeta}$.\footnote{Justified by the auxiliary claim
    \leanid{clm:g-comm-stability} in the cited paper.}
  \item Insert Bob's second measurement via $\closenessOfIP$, at cost
    $2\sqrt{\zeta}$.
  \item Apply point commutativity, at cost
    $6\sqrt{\gamma(m+1)}$.\footnote{The point-commutativity theorem,
    labelled \leanid{thm:commutativity-points}.}
  \item Remove the right-register $G^{\bx}$ total, at cost
    $\sqrt{\zeta} + 6\sqrt{\gamma(m+1)}$.\footnote{Justified by
    \leanid{clm:g-comm-stability2}.}
  \item Reverse the $\closenessOfIP$ insertion of $A$ on the left
    register, at cost $4\sqrt{\zeta}$ (steps~6--7 of the paper).
  \item Post-process $G$ measurement outcomes, at cost
    $\sqrt{\zeta}$.\footnote{Justified by an auxiliary fact derived
    in the formalization.}
  \item Reverse the second-coordinate post-processing, at cost
    $\sqrt{\zeta}$.
\end{enumerate}
The individual error terms sum to
$12\sqrt{\zeta} + 12\sqrt{\gamma(m+1)}$. Multiplying by $2$ (from
the commutator-square expansion) and simplifying $12\sqrt{m+1} \leq
24m$ gives $\nu = 48m\,(\sqrt{\gamma}+\sqrt{\zeta})$.

\section{The formal construction}

The formal lemma replicates the same chain,\footnote{The lemma
carries the name \leanid{commDataProcessedG\_of\_commutativityPoints} and is
implemented in
\path{MIPStarRE/LDT/Commutativity/ScalarApproximation/ProcessedG.lean}. The error parameter, defined in
\path{MIPStarRE/LDT/Commutativity/Scaffold/Core.lean} under the name
\leanid{commDataProcessedGError}, evaluates to
$48\,m\,(\gamma^{1/2} + \zeta^{1/2})$.}
each paper step recorded as a named phase block. The final assembly
adds these phases with the factor $2$ from the commutator-square
expansion, and the total matches the paper exactly.

The one substep that deserves explicit mention is the transition
from the $\BAB$ term to the $\ABA$ term. In the paper this is
justified by the projectivity of $G$, which makes
\[
  \E_{\bu,\bv,\bx,\by} \sum_{a,b} G^{\bx}_{[g(\bu)=a]}
  G^{\by}_{[h(\bv)=b]}
  = 
  \E_{\bu,\bv,\bx,\by} \sum_{a,b} G^{\by}_{[h(\bv)=b]}
  G^{\bx}_{[g(\bu)=a]},
\]
an exact equality. The formal proof uses the corresponding swap
lemma,\footnote{Its declaration is
\leanid{evaluatedSliceCommutation\_avg\_swap\_terms} in
\path{MIPStarRE/LDT/Commutativity/EvaluatedSliceCommutation/Averages.lean},
line~217.}
which proves the average of the $\BAB$ term and the average of the
$\ABA$ term equal as an \emph{equality} (not an approximation),
exactly matching the paper's step.

\section{Conclusion}

There is no mathematical divergence. The formal chain follows the
paper's proof step-by-step, uses the same auxiliary lemmas (closeness
of inner products, point commutativity, and $G$-commutativity
stability), and achieves the same total error bound
$48m\,(\sqrt{\gamma}+\sqrt{\zeta})$. The swap identity between the
$\BAB$ and $\ABA$ averages is an exact equality in both the paper
and the formalization. This note therefore documents a successful
alignment, not a gap.

\bibliographystyle{alpha}
\bibliography{references}

\end{document}